\chapter{Review of the Literature}\label{chap:literature_review}

% Creates a chapter and allows referencing with \autoref{chap:literature_review}

\section{Example Section}

Lorem ipsum dolor sit amet, consectetur adipiscing elit.\supercite{example_reference_1}

% \supercite{} inserts a citation from your .bib file

The major components of this framework are summarized in \textbf{\autoref{tab:example_table}}.

% \autoref{} automatically formats the reference as "Table X"
% Always reference tables/figures in text before they appear

\begin{table}[H]
\setstretch{1.0} % Controls vertical spacing within the table
\centering % Centers the table on the page

\caption{Example Table Title for Demonstration Purposes. This caption appears in the List of Tables}
% Caption appears above the table (standard for dissertations)

\label{tab:example_table}
% Label MUST come after caption for correct cross-referencing

\begin{tabular}{p{3.1cm} p{3.5cm} p{6.55cm}}
% Defines column layout:
% p{} creates fixed-width columns with automatic text wrapping
% This is preferred over l/c/r for long text in dissertations

\hline
\textbf{Category} & \textbf{Components} & \textbf{Description} \\
\hline

Example Row 1 
& Example Elements 
& Lorem ipsum dolor sit amet, consectetur adipiscing elit.\supercite{example_reference_1} 
\\

Example Row 2 
& Example Elements 
& Sed do eiusmod tempor incididunt ut labore et dolore magna aliqua.\supercite{example_reference_2} 
\\

Example Row 3 
& Example Elements 
& Ut enim ad minim veniam, quis nostrud exercitation ullamco laboris.\supercite{example_reference_3} 
\\

\hline
\hline

\multicolumn{3}{p{.95\linewidth}}{
\textit{Note.} This row spans all columns using \texttt{\textbackslash multicolumn}. 
Use this section for abbreviations, clarifications, or additional notes about the table.
}
\\

\hline
\end{tabular}
\end{table}

% Key table concepts:
% - Columns are separated by &
% - Rows end with \\
% - \hline creates horizontal rules
% - \multicolumn merges columns across the table width
% - p{} columns prevent text overflow and improve readability

Lorem ipsum dolor sit amet, consectetur adipiscing elit.

\section{Another Section}

An example glossary term is \gls{HMSC}.

% \gls{} pulls from the glossary file (e.g., common/acronyms.tex)
% First use = full term, subsequent uses = abbreviation

Example inline math: $p < 0.05$

% Inline math is used within a sentence

Example display math:

\[
y = mx + b
\]

% Display math is centered and used for emphasis or equations

\section{Summary}

Lorem ipsum dolor sit amet, consectetur adipiscing elit.